%% ----------------------------------------------------------------------------
%%  §1.6  Funzioni invertibili. Funzioni monotòne
%% ----------------------------------------------------------------------------
\section{Funzioni invertibili. Funzioni monotòne}
\label{sec:invertibili}

\begin{marginfigure}
\centering
\begin{tikzpicture}[scale=0.62, every node/.style={font=\scriptsize}]
  \foreach \k/\tit in {0/{iniettiva, non suriettiva}, 1/{suriettiva, non iniettiva}, 2/{biunivoca}} {
  \begin{scope}[shift={(0,-\k*3.6)}]
    \draw[blu, thick] (0,0) ellipse (0.7 and 1.3);
    \draw[rame, thick] (3,0) ellipse (0.7 and 1.3);
    \node at (1.5,-1.75) {\tit};
    \ifnum\k=0
      \foreach \y in {0.7,0,-0.7} \fill[blu] (0,\y) circle (1.4pt);
      \foreach \y in {0.9,0.3,-0.3,-0.9} \fill[rame] (3,\y) circle (1.4pt);
      \draw[-Latex,grigio] (0.1,0.7) -- (2.9,0.9);
      \draw[-Latex,grigio] (0.1,0) -- (2.9,-0.3);
      \draw[-Latex,grigio] (0.1,-0.7) -- (2.9,-0.9);
    \fi
    \ifnum\k=1
      \foreach \y in {0.9,0.3,-0.3,-0.9} \fill[blu] (0,\y) circle (1.4pt);
      \foreach \y in {0.7,0,-0.7} \fill[rame] (3,\y) circle (1.4pt);
      \draw[-Latex,grigio] (0.1,0.9) -- (2.9,0.7);
      \draw[-Latex,grigio] (0.1,0.3) -- (2.9,0.7);
      \draw[-Latex,grigio] (0.1,-0.3) -- (2.9,0);
      \draw[-Latex,grigio] (0.1,-0.9) -- (2.9,-0.7);
    \fi
    \ifnum\k=2
      \foreach \y in {0.7,0,-0.7} \fill[blu] (0,\y) circle (1.4pt);
      \foreach \y in {0.7,0,-0.7} \fill[rame] (3,\y) circle (1.4pt);
      \draw[-Latex,grigio] (0.1,0.7) -- (2.9,0);
      \draw[-Latex,grigio] (0.1,0) -- (2.9,-0.7);
      \draw[-Latex,grigio] (0.1,-0.7) -- (2.9,0.7);
    \fi
  \end{scope}}
\end{tikzpicture}
\caption{Iniettiva: frecce che non si incontrano mai. Suriettiva: ogni elemento di $B$
è raggiunto. Biunivoca: entrambe le cose, e le frecce si possono «invertire».}
\label{fig:iniettiva-suriettiva}
\end{marginfigure}

\begin{definizione}[Funzioni iniettive, suriettive, biunivoche]\index{funzione!iniettiva}\index{funzione!suriettiva}\index{funzione!biunivoca}
Sia $f\colon A\to B$ una funzione.
\begin{itemize}
  \item $f$ si dice \emph{iniettiva} se elementi distinti hanno immagini distinte, cioè,
  equivalentemente, se per ogni $x_{1},x_{2}\in A$ sussiste l'implicazione
  \begin{equation}\label{eq:iniettiva}
    f(x_{1})=f(x_{2}) \ \implica\ x_{1}=x_{2}.
  \end{equation}
  \item $f$ si dice \emph{suriettiva} (da $A$ \emph{verso} $B$) se per ogni $y\in B$
  esiste almeno un $x\in A$ tale che $y=f(x)$.
  \item $f$ si dice \emph{biunivoca} (o \emph{biiettiva}) se è contemporaneamente
  iniettiva e suriettiva. Ciò vuol dire che $f$ non solo fa corrispondere a ogni $x\in A$
  uno e un solo valore $y\in B$, ma anche che per ogni $y\in B$ esiste \emph{un solo}
  $x\in A$ tale che $y=f(x)$.
\end{itemize}
\end{definizione}


\begin{definizione}[Funzione inversa]\index{funzione!inversa}
Se $f\colon A\to B$ è biunivoca diciamo che $f$ è \emph{invertibile}. La funzione da $B$
ad $A$ che a ogni $y\in B$ fa corrispondere l'unico $x\in A$ per cui $f(x)=y$ si chiama
\emph{funzione inversa} e si indica con $f^{-1}$. Si ha (il simbolo $\perogni$ si legge
\emph{per ogni})
\begin{equation}\label{eq:inversa}
  f^{-1}(f(x))=x \quad\perogni x\in A; \qquad\qquad f(f^{-1}(y))=y \quad\perogni y\in B.
\end{equation}
\end{definizione}

Da notare che spesso si cambiano le notazioni e, invece di usare $x=f^{-1}(y)$, si
preferisce --- con un puro scambio di simboli --- scrivere $y=f^{-1}(x)$, per mettere in
luce che la variabile indipendente è quella dell'argomento di $f^{-1}$. Con questa
convenzione il grafico di $f^{-1}$ si ottiene da quello di $f$ scambiando il ruolo degli
assi, cioè con una simmetria rispetto alla bisettrice $y=x$ (\cref{fig:inversa}).

\begin{figure}[htb]
\centering
\begin{tikzpicture}
\begin{axis}[width=6.2cm, height=6.2cm, xmin=-0.3, xmax=3.2, ymin=-0.3, ymax=3.2,
  xtick=\empty, ytick=\empty, title={\footnotesize $y=f(x)$}]
  \addplot[tratteggio, domain=-0.3:3.2] {x};
  \addplot[curva, domain=0.5:2.5, samples=80] {0.9+0.5*(x-0.5)+0.275*(x-0.5)^2};
  \addplot[tratteggio] coordinates {(0.5,0) (0.5,0.9) (0,0.9)};
  \addplot[tratteggio] coordinates {(2.5,0) (2.5,3.0) (0,3.0)};
  \draw[blu, thick, decorate, decoration={brace, amplitude=3pt, mirror}] (axis cs:0.5,-0.15) -- (axis cs:2.5,-0.15) node[midway, below=3pt, font=\scriptsize] {$A$};
  \draw[rame, thick, decorate, decoration={brace, amplitude=3pt}] (axis cs:-0.15,0.9) -- (axis cs:-0.15,3.0) node[midway, left=3pt, font=\scriptsize] {$B$};
\end{axis}
\end{tikzpicture}\hfill
\begin{tikzpicture}
\begin{axis}[width=6.2cm, height=6.2cm, xmin=-0.3, xmax=3.2, ymin=-0.3, ymax=3.2,
  xtick=\empty, ytick=\empty, title={\footnotesize $y=f^{-1}(x)$}]
  \addplot[tratteggio, domain=-0.3:3.2] {x};
  % inversa: scambio delle coordinate, disegnata in forma parametrica
  \addplot[curva2, domain=0.5:2.5, samples=80] ({0.9+0.5*(x-0.5)+0.275*(x-0.5)^2},{x});
  \addplot[tratteggio] coordinates {(0.9,0) (0.9,0.5) (0,0.5)};
  \addplot[tratteggio] coordinates {(3.0,0) (3.0,2.5) (0,2.5)};
  \draw[rame, thick, decorate, decoration={brace, amplitude=3pt, mirror}] (axis cs:0.9,-0.15) -- (axis cs:3.0,-0.15) node[midway, below=3pt, font=\scriptsize] {$B$};
  \draw[blu, thick, decorate, decoration={brace, amplitude=3pt}] (axis cs:-0.15,0.5) -- (axis cs:-0.15,2.5) node[midway, left=3pt, font=\scriptsize] {$A$};
\end{axis}
\end{tikzpicture}
\caption{Il grafico di una funzione invertibile $f\colon A\to B$ e quello della sua
inversa $f^{-1}\colon B\to A$: l'uno è il simmetrico dell'altro rispetto alla
bisettrice del primo e terzo quadrante (tratteggiata).}
\label{fig:inversa}
\end{figure}

\begin{esempio}[Le funzioni \eqref{eq:f1}--\eqref{eq:f4} sono invertibili?]
La funzione \eqref{eq:f1}, $f(x)=2x+1$, è invertibile da $\R$ in $\R$: fissato
$y=f(x)$, risulta $x=(y-1)/2$, e questo $x$ è unico. La funzione inversa è quindi
\begin{equation}\label{eq:inversa-f1}
  f^{-1}(x)=\frac{x-1}{2}.
\end{equation}
Anche la funzione \eqref{eq:f2}, $f(x)=1/x$, è invertibile da $\R\setminus\{0\}$ in
$\R\setminus\{0\}$ e risulta $f^{-1}(x)=1/x$: è solo un caso che $f$ e $f^{-1}$
coincidano! La funzione \eqref{eq:f3} è invertibile da $[0,+\infty)$ in $[0,+\infty)$
e $f^{-1}(x)=x^{2}$, per $x\ge0$. La \eqref{eq:f4} non è invertibile, perché non
stabilisce una corrispondenza biunivoca tra $\R$ e l'insieme $B=\{0,1\}$ dei valori:
tutti gli interi, per esempio, hanno la stessa immagine $0$.
\end{esempio}

\begin{attenzione}[$f^{-1}$ non è $1/f$]
Il simbolo $f^{-1}$ indica la funzione inversa e \emph{non} il reciproco $1/f$. Per
$f(x)=2x+1$ si ha $f^{-1}(x)=(x-1)/2$, mentre $1/f(x)=1/(2x+1)$: funzioni del tutto
diverse. La notazione è infelice ma universale, e vale la pena di abituarsi subito.
\end{attenzione}

\begin{notastorica}[Cantor: tanti punti in un segmento quanti in un quadrato]
Le corrispondenze biunivoche sono lo strumento con cui Georg Cantor, a partire dal
1874, imparò a confrontare gli insiemi infiniti: due insiemi hanno la stessa
\emph{potenza} --- per gli insiemi finiti, lo stesso numero di elementi --- quando
esiste una funzione biunivoca dall'uno all'altro. Nel 1877 Cantor si chiese se il
quadrato $[0,1]\times[0,1]$ abbia «più punti» del segmento $[0,1]$, come tutti davano
per scontato, e trovò invece una corrispondenza biunivoca tra i due. L'idea, nella forma
più semplice, è di intrecciare le cifre decimali: al punto
$(0{,}a_{1}a_{2}a_{3}\dots,\ 0{,}b_{1}b_{2}b_{3}\dots)$ si associa il numero
$0{,}a_{1}b_{1}a_{2}b_{2}a_{3}b_{3}\dots$. Dedekind, a cui Cantor aveva sottoposto la
dimostrazione, osservò che, per via dei numeri con due scritture decimali (come
$0{,}1=0{,}0999\dots$), una corrispondenza costruita così è iniettiva ma non suriettiva;
Cantor rimediò con un argomento più elaborato e il 29 giugno 1877 gli scrisse: «Je le
vois, mais je ne le crois pas!» («Lo vedo, ma non ci credo!»). Il risultato, pubblicato
nel 1878, sembrava togliere significato alla nozione di dimensione; la questione fu
chiarita solo nel 1911, quando L.~E.~J. Brouwer dimostrò che tra spazi di dimensione
diversa non esistono corrispondenze biunivoche continue insieme con la loro inversa.
\end{notastorica}

\subsection{Funzioni monotòne}

\begin{definizione}[Funzioni monotòne]\index{funzione!monotòna}\index{funzione!crescente}\index{funzione!decrescente}
Diremo che una funzione $f$ è \emph{monotòna} in un insieme $A$ se verifica una delle
condizioni seguenti, per ogni $x_{1},x_{2}\in A$:
\begin{align}
  &f \text{ \emph{strettamente crescente}:}   & x_{1}<x_{2} &\implica f(x_{1})<f(x_{2}),\label{eq:str-cresc}\\
  &f \text{ \emph{crescente}:}                & x_{1}<x_{2} &\implica f(x_{1})\le f(x_{2}),\label{eq:cresc}\\
  &f \text{ \emph{strettamente decrescente}:} & x_{1}<x_{2} &\implica f(x_{1})>f(x_{2}),\label{eq:str-decr}\\
  &f \text{ \emph{decrescente}:}              & x_{1}<x_{2} &\implica f(x_{1})\ge f(x_{2}).\label{eq:decr}
\end{align}
Una funzione che verifica la \eqref{eq:str-cresc} oppure la \eqref{eq:str-decr} si
dice \emph{strettamente monotòna}.
\end{definizione}

\begin{marginfigure}
\centering
\begin{tikzpicture}
\begin{axis}[width=5cm, height=4cm, xmin=0, xmax=4, ymin=0, ymax=3, xtick=\empty, ytick=\empty]
  \addplot[curva, domain=0.2:3.8, samples=60] {0.5+0.9*(x-0.2)^0.6};
  \addplot[curva2] coordinates {(0.2,0.6) (1,1.4) (2,1.4) (3,2.2) (3.8,2.2)};
\end{axis}
\end{tikzpicture}
\caption{In blu una funzione strettamente crescente; in rame una funzione crescente
ma non strettamente: ha dei «pianerottoli» su cui $f(x_1)=f(x_2)$ con $x_1<x_2$.}
\end{marginfigure}

Ad esempio, la funzione \eqref{eq:f1} è strettamente crescente su $\R$; la funzione
\eqref{eq:f2} è strettamente decrescente, \emph{separatamente} negli insiemi $\{x>0\}$ e
$\{x<0\}$ (ma non su tutto il suo dominio: $-1<1$ e $f(-1)=-1<1=f(1)$). La funzione
\eqref{eq:f3} è strettamente crescente. La funzione \eqref{eq:f4} non è monotòna su
$\R$. Una funzione costante è contemporaneamente crescente e decrescente, ma non
strettamente.

Ci sarà utile nel seguito un \emph{criterio} per riconoscere se una data funzione è
invertibile. Il risultato che segue è semplice, ma è uno dei più usati del libro: è
grazie a esso che potremo definire le radici, i logaritmi e le funzioni trigonometriche
inverse.

\begin{teorema}[Criterio di invertibilità]\label{teo:monotona-invertibile}\index{criterio di invertibilità}
Sia $f\colon A\to B$ una funzione suriettiva. Se $f$ è strettamente monotòna in $A$,
allora $f$ è invertibile, e l'inversa $f^{-1}\colon B\to A$ è strettamente monotòna
dello stesso tipo (crescente se $f$ è crescente, decrescente se $f$ è decrescente).
\end{teorema}

\begin{schema}
\Ipotesi{$f\colon A\to B$ suriettiva e strettamente monotòna.}
\Tesi{$f$ è iniettiva (dunque biunivoca, dunque invertibile) e $f^{-1}$ è strettamente
monotòna dello stesso tipo.}
\end{schema}

\begin{strategia}
Poiché $f$ è suriettiva per ipotesi, per avere l'invertibilità manca solo
l'iniettività: dobbiamo mostrare che due punti distinti $x_{1}\neq x_{2}$ hanno immagini
distinte. Ma se $x_{1}\neq x_{2}$, uno dei due è più piccolo, diciamo $x_{1}<x_{2}$
(dicotomia); e allora la stretta monotonia dice che $f(x_{1})<f(x_{2})$ oppure
$f(x_{1})>f(x_{2})$: in ogni caso $f(x_{1})\neq f(x_{2})$. Per la monotonia
dell'inversa ragioniamo per assurdo: se $y_{1}<y_{2}$ ma $f^{-1}(y_{1})\ge f^{-1}(y_{2})$,
applicando $f$ (crescente) otterremmo $y_{1}\ge y_{2}$.
\end{strategia}

\begin{dimostrazione}
Supponiamo, per fissare le idee, che $f$ sia strettamente crescente (il caso
decrescente è identico, con le disuguaglianze invertite). Siano $x_{1},x_{2}\in A$ con
$x_{1}\neq x_{2}$. Per la dicotomia \ref{ass:dicotomia} è $x_{1}<x_{2}$ oppure
$x_{2}<x_{1}$; nel primo caso $f(x_{1})<f(x_{2})$, nel secondo $f(x_{2})<f(x_{1})$. In
entrambi i casi $f(x_{1})\neq f(x_{2})$: $f$ è iniettiva. Essendo anche suriettiva, è
biunivoca e quindi invertibile.

Proviamo che $f^{-1}$ è strettamente crescente. Siano $y_{1},y_{2}\in B$ con
$y_{1}<y_{2}$ e poniamo $x_{1}=f^{-1}(y_{1})$, $x_{2}=f^{-1}(y_{2})$, così che
$f(x_{1})=y_{1}$ e $f(x_{2})=y_{2}$. Se fosse $x_{1}\ge x_{2}$, allora per la stretta
crescenza di $f$ (e banalmente, nel caso $x_{1}=x_{2}$) sarebbe $f(x_{1})\ge
f(x_{2})$, cioè $y_{1}\ge y_{2}$, contro l'ipotesi. Dunque $x_{1}<x_{2}$, cioè
$f^{-1}(y_{1})<f^{-1}(y_{2})$.
\end{dimostrazione}

\begin{attenzione}[Il viceversa è falso]
Una funzione invertibile non è necessariamente monotòna. La funzione
$f\colon[0,2]\to[0,2]$ definita da $f(x)=x$ per $x\in[0,1)$ e $f(x)=3-x$ per
$x\in[1,2]$ è biunivoca (il lettore ne disegni il grafico) ma non è monotòna. Per le
funzioni \emph{continue su un intervallo}, invece, invertibilità e stretta monotonia sono
equivalenti: nel capitolo sulle funzioni continue dimostreremo che una funzione continua
e strettamente monotòna su un intervallo è invertibile (\cref{teo:invertibilita}); anche
il viceversa è vero, come conseguenza del teorema dei valori intermedi.
\end{attenzione}

\subsection{Immagini, immagini inverse, funzioni composte}

Concludiamo il paragrafo con alcune utili definizioni e notazioni.

\begin{definizione}[Immagine e immagine inversa di un insieme]\index{immagine!di un insieme}\index{codominio}\index{immagine inversa}
Sia $f\colon A\to B$ una funzione. Se $X$ è un sottoinsieme di $A$, l'\emph{immagine}
di $X$ mediante $f$, indicata con $f(X)$, è il sottoinsieme di $B$ definito da
\begin{equation}\label{eq:immagine}
  f(X)=\{y\in B : \esiste x\in X : y=f(x)\}
\end{equation}
(il simbolo $\esiste$ si legge \emph{esiste}). L'insieme $f(A)$, cioè l'immagine di $A$
mediante $f$, si chiama \emph{codominio} o \emph{immagine} di $f$. Evidentemente la
funzione $f$ è suriettiva da $A$ verso $B$ se e solo se $f(A)=B$. Se $Y$ è un
sottoinsieme di $B$, l'\emph{immagine inversa} (o \emph{controimmagine}) di $Y$ mediante
$f$, indicata con $f^{-1}(Y)$, è il sottoinsieme di $A$ definito da
\begin{equation}\label{eq:controimmagine}
  f^{-1}(Y)=\{x\in A : f(x)\in Y\}.
\end{equation}
\end{definizione}

\begin{osservazione}
La notazione $f^{-1}(Y)$ ha senso \emph{anche se $f$ non è invertibile}: $f^{-1}(Y)$ è
un insieme, non il valore di una funzione. Ad esempio, per $f(x)=x^{2}$ si ha
$f^{-1}(\{4\})=\{-2,2\}$ e $f^{-1}(\{-1\})=\emptyset$. Allo stesso modo, nel linguaggio
degli ingegneri, l'immagine inversa di un intervallo di valori «accettabili» in uscita è
l'insieme degli ingressi che producono quelle uscite, indipendentemente dal fatto che la
legge ingresso--uscita sia invertibile.
\end{osservazione}

\begin{definizione}[Funzione composta]\index{funzione!composta}
Siano $X,Y,Z$ tre insiemi e siano $g\colon X\to Y$ e $f\colon Y\to Z$ due funzioni, tali
che l'insieme $Y$ contenente i valori della prima coincida con il dominio della seconda.
Allora si può considerare la \emph{funzione composta} $h\colon X\to Z$, definita da
$h(x)=f(g(x))$ per $x\in X$. In tal caso si usa la notazione $h=f\circ g$; in altre
parole, si pone
\[
  (f\circ g)(x)=f(g(x)) \qquad\perogni x\in X.
\]
\end{definizione}

\begin{marginfigure}
\centering
\begin{tikzpicture}[scale=0.8, every node/.style={font=\scriptsize}]
  \node[draw, blu, rounded corners=2pt, minimum width=1cm] (X) at (0,0) {$X$};
  \node[draw, blu, rounded corners=2pt, minimum width=1cm] (Y) at (2,0) {$Y$};
  \node[draw, blu, rounded corners=2pt, minimum width=1cm] (Z) at (4,0) {$Z$};
  \draw[-Latex] (X) -- node[above] {$g$} (Y);
  \draw[-Latex] (Y) -- node[above] {$f$} (Z);
  \draw[-Latex, rame] (X) to[bend right=40] node[below] {$f\circ g$} (Z);
\end{tikzpicture}
\caption{La composizione $f\circ g$: prima si applica $g$, poi $f$. Si legge «$f$
composto $g$» e si calcola da destra verso sinistra.}
\end{marginfigure}

\begin{esempio}
Siano $g(x)=2x+1$ e $f(y)=\sqrt{y}$. Allora $(f\circ g)(x)=\sqrt{2x+1}$, definita per
$x\ge-\tfrac12$, mentre $(g\circ f)(x)=2\sqrt{x}+1$, definita per $x\ge0$: la
composizione \emph{non è commutativa}. Con la notazione della composizione, le
relazioni \eqref{eq:inversa} che caratterizzano la funzione inversa si scrivono in modo
compatto come $f^{-1}\circ f=\mathrm{id}_{A}$ e $f\circ f^{-1}=\mathrm{id}_{B}$, dove
$\mathrm{id}$ indica la funzione \emph{identità} $x\mapsto x$.
\end{esempio}

\begin{python}[Comporre e invertire funzioni con \texttt{sympy}]
# sympy manipola le espressioni simbolicamente, come faremmo a mano.
import sympy as sp
x, y = sp.symbols("x y", real=True)

g = 2*x + 1
f = sp.sqrt(x)
print("f(g(x)) =", f.subs(x, g))         # composizione f o g
print("g(f(x)) =", g.subs(x, f))         # composizione g o f

# Inversa di g: risolviamo y = g(x) rispetto a x
inversa = sp.solve(sp.Eq(y, g), x)[0]
print("g^-1(y) =", inversa)
# Verifica delle relazioni g^-1(g(x)) = x e g(g^-1(y)) = y
print(sp.simplify(inversa.subs(y, g)),
      sp.simplify(g.subs(x, inversa)))
\end{python}
\begin{risultato}
f(g(x)) = sqrt(2*x + 1)
g(f(x)) = 2*sqrt(x) + 1
g^-1(y) = y/2 - 1/2
x y
\end{risultato}

\begin{esercizio}
Sia $f\colon\R\to\R$, $f(x)=x^{3}$. Dimostrare, usando solo la definizione, che $f$ è
strettamente crescente, e dedurne che è iniettiva. (La suriettività, cioè l'esistenza
della radice cubica di ogni reale, richiede l'estremo superiore e sarà provata nel
\cref{teo:radice-n}.)
\begin{soluzione}
Siano $x_{1}<x_{2}$. Allora $x_{2}-x_{1}>0$ e
$x_{2}^{3}-x_{1}^{3}=(x_{2}-x_{1})(x_{2}^{2}+x_{1}x_{2}+x_{1}^{2})$. Il secondo fattore
è $\bigl(x_{2}+\tfrac12x_{1}\bigr)^{2}+\tfrac34x_{1}^{2}\ge0$, e si annulla solo se
$x_{1}=x_{2}=0$, escluso. Dunque il prodotto è $>0$, cioè $x_{1}^{3}<x_{2}^{3}$.
L'iniettività segue dalla stretta crescenza, come nella prima parte della dimostrazione
del \cref{teo:monotona-invertibile}: se $x_{1}\neq x_{2}$, uno dei due è minore
dell'altro, e le immagini sono diverse. \emph{Scelta logica:} per confrontare
$x_{1}^{3}$ e $x_{2}^{3}$ si studia il segno della differenza
(\cref{prop:ordine-differenza}); la fattorizzazione di $x_{2}^{3}-x_{1}^{3}$ riduce il
problema al segno di un trinomio, che il completamento del quadrato rende evidente.
\end{soluzione}
\end{esercizio}
